\section{Extracting multiplicatively organized pairs}\label{sec:extraction}

We now use \Cref{com:variation,com:rare-positive}: the ordinary center-star
variation satisfies \(C_Q=o(S)\), whereas the point exception mass satisfies
\(\liminf \E_\mu H_p^*/S>0\).  Here \(S\to\infty\), and the parameters,
levels, and exception indicators are those of \Cref{sec:budgets}.
Our objective is a graph on two copies of \(P\) whose edges carry ratios of
unit jumps.  These ratios will have large height, while their alternating
products around every complete rectangle have negligible height.

Throughout this section an \emph{ordinary wedge} is sampled by taking
\(q\) with law \(\kappa\) and then taking \(p,r\) independently and uniformly
among its neighbors in \(E\).  Its two point marginals are \(\mu\).
We retain the notation \(W_c\) for the width of common levels, \(W_b\) for
the width of damaged rare levels, and \(m_{pq}\) for the integrated
inconsistency indicator of the edge \(pq\).  In particular,
\begin{equation}\label{ext:inputs}
 \E_\mu H_p^*+\E_\kappa H_q^*=O(S),\qquad
 \E_E m_{pq}=O(1),\qquad
 W_c\le S/\epsilon,\qquad W_b=o(S),
 \quad \epsilon=k^{-1/16}.
\end{equation}
All integrals over levels include the prescribed place weights.

\subsection{A pointwise estimate and its tail consequence}

The first-moment bound in \eqref{ext:inputs} alone does not prevent all
exception mass from being carried by extremely rare points.  We obtain the
needed tail control by applying the product formula separately to each
distinct pair of neighbors of a point.

\begin{lemma}[Pointwise rare-level estimate]\label{ext:pointwise}
For every \(p\in P\) and every pair of distinct neighbors \(q,q'\in Q\),
\begin{equation}\label{ext:pointwise-bound}
 2H_p^*\le
 3\bigl(H_q^*+H_{q'}^*+m_{pq}+m_{pq'}\bigr)
       +W_c+W_b+2\log 2.
\end{equation}
\end{lemma}

\begin{proof}
At a good rare level, let \(X,U,V\) be the exception indicators of
\(p,q,q'\), and let \(I,I'\) be the inconsistency indicators of the two
edges.  Write \(P_+\) for the rare positive charge in the point-star score:
it equals two when \(p\) is exceptional, both bits are minority bits, and
the two neighbors have matching nonshared coordinates; otherwise it is
zero.  Write \(N=\abs{Y_{pq}-Y_{pq'}}\).
If \(U=V=I=I'=0\), the two bits are the same function of \(X\): they are
both majority bits for \(X=0\), and both minority bits for \(X=1\).
In the latter case the nonshared coordinate of both neighbors has the
dominant label, so \(P_+=2\).  Consequently, pointwise,
\[
 P_+\ge 2X-2(U+V+I+I'),\qquad
 N\le U+V+I+I',
\]
and hence
\begin{equation}\label{ext:local-rare-score}
 P_+-N\ge 2X-3(U+V+I+I').
\end{equation}
In particular, a simultaneous exception at \(p\) and at a neighbor is
charged directly to the neighbor indicator.

For completeness, the upper bound for this rare-only score is also
pointwise.  Put \(d=z_p-z_q\), \(d'=z_p-z_{q'}\).  These jumps are distinct
and nonzero.  At a place \(v\), sort their thresholds as \(b\le b'\), and set
\[
 u=-\log\abs{d-d'}_v,\qquad \alpha=u-b,\qquad
 T_v=\begin{cases}\alpha+\log 2,&v\text{ archimedean},\\
                    \alpha,&v\text{ nonarchimedean}.
       \end{cases}
\]
The triangle inequality makes \(T_v\ge0\).  Matching coordinates in a
square cell of side length \(r\) gives distance at most \(\sqrt2 r\), and
matching nonarchimedean labels gives distance at most \(r\).
Thus positive charges with both bits one lie within \([b-T_v,b]\), and
positive charges with both bits zero lie within \([b',b'+T_v]\).

Only one of these two intervals can contain a rare positive charge.
Indeed every level where zero is the minority bit precedes every level
where one is the minority bit.  A zero-bit rare positive would require a
threshold at least \(b'\), while a one-bit rare positive would require a
later threshold strictly less than \(b\), which is impossible.
The total rare positive charge at this place is therefore at most
\(2T_v\).  The all-level negative charge is \(b'-b\), because the level
interval contains both thresholds.  Moreover,
\[
 -\log\abs{(d-d')(1/d-1/d')}_v=2\alpha-(b'-b).
\]
Summing over places, the product formula cancels this last expression.
The archimedean weights sum to one, so the rare positives minus the
all-level negative charge total at most \(2\log2\).

Integrating \eqref{ext:local-rare-score} gives a lower bound for that same
score.  Outside the good rare levels there are no positive terms, and
the negative charge is at most \(W_c+W_b\).  Comparing the lower and upper
bounds proves \eqref{ext:pointwise-bound}.
\end{proof}

\begin{lemma}[Tail control]\label{ext:tail}
Let \(R_0=\epsilon^{-2}=k^{1/8}\).  Then
\begin{equation}\label{ext:tail-bound}
 \E_\mu\bigl[H_p^*\mathbf1_{\{H_p^*>R_0S\}}\bigr]=o(S).
\end{equation}
\end{lemma}

\begin{proof}
Choose an absolute \(C_0\) for which each of the two exception expectations
in \eqref{ext:inputs} is at most \(C_0S\), for all sufficiently large
indices.  Put
\[
 T=\{p:H_p^*>R_0S\},\qquad
 A_t=\E_\mu[H_p^*\mathbf1_T],\qquad
 d_{P,\min}=\min_{p\in P}d(p).
\]
On \(T\), the deterministic last three terms of
\eqref{ext:pointwise-bound} have sum at most \(H_p^*\), eventually, since
\[
 \frac{W_c+W_b+2\log2}{R_0S}
 \le\epsilon+\frac{W_b+2\log2}{R_0S}=o(1).
\]
After absorbing them, we may cap each neighbor height at \(H_p^*\):
\begin{equation}\label{ext:capped-pointwise}
 H_p^*\le3\bigl(\min(H_p^*,H_q^*)+\min(H_p^*,H_{q'}^*)
                       +m_{pq}+m_{pq'}\bigr),\qquad p\in T,
\end{equation}
for distinct \(q,q'\).  If a cap changes a term, that capped term already
equals \(H_p^*\) and suffices for this inequality; otherwise the preceding
inequality is unchanged.

Sample \(p\) with law \(\mu\) and two neighbors independently and uniformly
given \(p\).  Restrict the left side to distinct neighbors.  The exact
conditional probability of distinctness is \(1-1/d(p)\).  On the
nonnegative right side we can drop this restriction.  Thus
\begin{equation}\label{ext:tail-average}
 (1-d_{P,\min}^{-1})A_t
 \le 6\E_E[\min(H_p^*,H_q^*)\mathbf1_T]+C_m,
\end{equation}
where \(C_m\) is an absolute upper bound for \(6\E_E m_{pq}\).
This calculation uses the distinctness probability at each individual
point, so it remains valid even when its exception mass is very large.

We estimate the capped expectation by layer cake.  For \(u\ge0\) let
\[
 F_u=\{p:H_p^*>\max(u,R_0S)\},\qquad
 G_u=\{q:H_q^*>u\},
\]
and write \(h_P(u)=\mu(F_u)\), \(h_Q(u)=\kappa(G_u)\).
The marginal version of the expansion estimate from \Cref{sec:incidence},
with \(\eta=\gamma-1/2>0\), is
\[
 \Prob_E(F,G)\le\sqrt{\mu(F)\kappa(G)}
 \left[\gamma^{-(1+2\eta)}(\mu(F)\kappa(G))^\eta
                         +\gamma^{-1}\xi_t\right].
\]
Because \(h_P(u)\le C_0/R_0\) and \(h_Q(u)\le1\), the same multiplier
\[
 \vartheta_t=\gamma^{-(1+2\eta)}(C_0/R_0)^\eta
                      +\gamma^{-1}\xi_t=o(1)
\]
works at every threshold \(u\), including thresholds where the center set
has full measure.  Tonelli's theorem and Cauchy--Schwarz now give
\begin{align*}
 \E_E[\min(H_p^*,H_q^*)\mathbf1_T]
 &=\int_0^\infty\Prob_E(F_u,G_u)\dd u\\
 &\le\vartheta_t
       \left(\int_0^\infty h_P(u)\dd u\right)^{1/2}
       \left(\int_0^\infty h_Q(u)\dd u\right)^{1/2}\\
 &=\vartheta_t\sqrt{A_t\E_\kappa H_q^*}
 \le\vartheta_t\sqrt{C_0A_tS}.
\end{align*}
The first integral equals the unknown tail first moment \(A_t\), including
its contribution below the threshold \(R_0S\).
Since \(d_{P,\min}\to\infty\), \eqref{ext:tail-average} implies eventually
\[
 \tfrac12 A_t\le6\vartheta_t\sqrt{C_0A_tS}+C_m
       \le\tfrac14 A_t+36C_0\vartheta_t^2S+C_m.
\]
The last step is Young's inequality.  Consequently
\[
 \frac{A_t}{S}\le144C_0\vartheta_t^2+\frac{4C_m}{S}\longrightarrow0,
\]
as required.
\end{proof}

\subsection{Point profiles and a common height band}

We next express a wedge's jump ratio as a difference of profiles attached
to its two point endpoints.  On a good rare level \((v,s)\), let
\(X_x(v,s)=I_x(v,s)\) be the exception indicator of the vertex \(x\)
from \Cref{sec:budgets}, and let
\(\sigma(v,s)=1\) when one is the minority bit and \(-1\) when zero is
the minority bit.  Define a real-valued profile, indexed by places, by
\begin{equation}\label{ext:profile-definition}
 a_p(v)=\int_{\substack{s:\,(v,s)\text{ is}\\\text{good rare}}}
                       \sigma(v,s)X_p(v,s)\dd s.
\end{equation}
For a place profile \(b\), write \(\norm{b}_1=\sum_v^*\abs{b(v)}\).

\begin{lemma}[Wedge profiles]\label{ext:profiles}
For an ordinary wedge \((p,q,r)\), there are nonnegative random variables
\(E_1,E_2\), with \(\E E_1+\E E_2=o(S)\), such that
\begin{align}
 \norm{(n_{pq}-n_{rq})-(a_p-a_r)}_1&\le E_1,
                                      \label{ext:profile-error}\\
 \norm{n_{pq}-n_{rq}}_1&\ge H_p^*+H_r^*-E_2.
                                      \label{ext:profile-size}
\end{align}
Here \(n_{pq}(v)=-\log\abs{z_p-z_q}_v\).
\end{lemma}

\begin{proof}
Let \(\mathcal G\) denote the good rare levels, and let \(I_{pq}\) be
the levelwise inconsistency indicator.  At a level in \(\mathcal G\),
write \(b\in\{0,1\}\) for the majority bit.  Consistency gives
\[
 Y_{pq}=b+\sigma(X_p+X_q)
\]
unless both endpoints are exceptional or the edge is inconsistent.
In all cases the absolute residual in this identity is at most
\(2(X_pX_q+I_{pq})\).  Subtracting the identities for the two edges
cancels the center contribution.  At each place,
\[
 n_{pq}(v)-n_{rq}(v)=\int(Y_{pq}-Y_{rq})\dd s,
\]
since both bits are single thresholds inside the level interval.
It follows that \eqref{ext:profile-error} holds with
\[
 E_1=\int_{\mathcal C}\abs{Y_{pq}-Y_{rq}}+W_b
       +2\int_{\mathcal G}(X_pX_q+X_rX_q+I_{pq}+I_{rq}).
\]
Its mean is \(o(S)\): the first term averages to \(C_Q=o(S)\), the
joint edge exceptions have integrated expectation \(o(S)\), and the
remaining bounds are in \eqref{ext:inputs}.

There is also the exact identity
\[
 \norm{n_{pq}-n_{rq}}_1=\int\abs{Y_{pq}-Y_{rq}}.
\]
On a good rare level with exactly one exceptional point endpoint, a
normal center, and two consistent edges, the bits differ.  More
explicitly, at every good rare level,
\[
 \abs{Y_{pq}-Y_{rq}}\ge X_p+X_r
 -2(X_pX_r+X_pX_q+X_rX_q+I_{pq}+I_{rq}).
\]
If none of the terms subtracted is present, this is the preceding
observation; if any is present, the right side is nonpositive.
Integrating proves \eqref{ext:profile-size} with
\[
 E_2=2\int_{\mathcal G}
                   (X_pX_r+X_pX_q+X_rX_q+I_{pq}+I_{rq}).
\]
The integrated joint point exceptions on a wedge and joint edge
exceptions are \(o(S)\), by \Cref{bud:rare-moments}.
Thus \(\E E_2=o(S)\) as well.
\end{proof}

The next selection makes these average error bounds uniform before
choosing a dyadic band.  This order matters because the probability of
the eventual band tends to zero.

\begin{lemma}[A band of accurate wedges]\label{ext:band}
There are a constant \(c>0\), independent of \(t\) along the fixed
sequence, a sequence \(\delta_t\to0\), and scales \(J\to\infty\) such
that a set of ordinary wedges of probability at least
\begin{equation}\label{ext:band-mass}
 \frac{c}{R_0\log(2R_0)}=t^{-o(1)}
\end{equation}
satisfies, uniformly on the set,
\begin{equation}\label{ext:band-conditions}
 J\le H_p^*+H_r^*\le2J,\qquad
 E_1+E_2\le\delta_t(H_p^*+H_r^*).
\end{equation}
Every such wedge has \(p\ne r\).
\end{lemma}

\begin{proof}
Set \(W=H_p^*+H_r^*\) and \(Z=E_1+E_2\).  By
\Cref{com:rare-positive}, choose \(c_0>0\), independent of \(t\) along
the fixed sequence, such that \(\E W\ge4c_0S\) eventually.
Choose \(\delta_t\to0\) sufficiently
slowly that \(\E Z/\delta_t=o(S)\).  Before selecting any band, discard
wedges with \(Z>\delta_tW\).  Their lost weight is at most
\[
 \E[W\mathbf1_{\{Z>\delta_tW\}}]\le\frac{\E Z}{\delta_t}=o(S).
\]
Also discard \(W<c_0S\), losing at most \(c_0S\).
If \(W>2R_0S\), its larger summand exceeds \(R_0S\) and is at least
\(W/2\).  Thus \Cref{ext:tail} gives
\[
 \E[W\mathbf1_{\{W>2R_0S\}}]
 \le2\E[H_p^*\mathbf1_{\{H_p^*>R_0S\}}]
    +2\E[H_r^*\mathbf1_{\{H_r^*>R_0S\}}]=o(S).
\]
The surviving weighted mass is at least \(2c_0S\), whereas every
surviving wedge has \(W\le2R_0S\).  Its probability is therefore at
least \(c_0/R_0\).

Partition \([c_0S,2R_0S]\) into at most
\(1+\lceil\log_2(2R_0/c_0)\rceil=O(\log(2R_0))\) dyadic bands,
using half-open endpoints except at the final endpoint.  One band has
the probability in \eqref{ext:band-mass}; if \(J\) is its lower endpoint,
then \(J\ge c_0S\to\infty\), and \eqref{ext:band-conditions} holds.
The parameter hierarchy gives \(k=t^{o(1)}\), so the lower bound in
\eqref{ext:band-mass} is indeed \(t^{-o(1)}\).
Finally, if \(p=r\), the left side of \eqref{ext:profile-size} is zero,
whereas its right side is at least \((1-\delta_t)W>0\) for all
sufficiently large indices.  Such a wedge cannot survive.
\end{proof}

\subsection{Removing common unit circles and choosing witnesses}

The algebraic argument will require that a retained pair does not lie
on a common unit circle containing many points of \(P\).  This exclusion
must concern all real unit circles, including those whose centers do
not belong to \(Q\).

\begin{lemma}[Popular-circle deletion]\label{ext:circles}
Call a real unit circle \emph{popular} if it contains at least
\(t^{9/10}\) points of \(P\).  The ordinary wedge probability of
\(p,r\) lying on a common popular circle is at most
\begin{equation}\label{ext:circle-cost}
 t^{-3/5+o(1)}+t^{-9/5+o(1)}+\frac{2}{\min_{q\in Q}d(q)}.
\end{equation}
This is \(o(1/(R_0\log(2R_0)))\).  In particular, removing such wedges
from the set in \Cref{ext:band} preserves the lower probability bound
\eqref{ext:band-mass}, after decreasing its constant.
\end{lemma}

\begin{proof}
Write \(M\) for the number of popular circles.  Every pair of distinct
real points lies on at most two unit circles.  Counting pairs on popular
circles therefore gives
\[
 M\binom{\lceil t^{9/10}\rceil}{2}
 \le2\binom{\abs P}{2},\qquad M\le t^{1/5+o(1)},
\]
using \(\abs P=t^{1+o(1)}\).
The set \(P_{\rm mult}\) of points on at least two popular circles has
size at most \(2\binom M2\le t^{2/5+o(1)}\), because distinct circles
intersect in at most two points.  The degree bounds give
\[
 \max_p\mu(p)\le t^{-1+o(1)},\qquad
 \max_q\kappa(q)\le t^{-2+o(1)}.
\]
Consequently wedges hitting \(P_{\rm mult}\) at either point endpoint
have probability at most \(t^{-3/5+o(1)}\).  Centers in \(Q\) whose
own unit circle is popular number at most \(M\); wedges with such a
center have probability at most \(t^{-9/5+o(1)}\).

To bound the remaining bad wedges, fix \(p\notin P_{\rm mult}\) and
a center \(q\) whose own circle is not popular, in the original sampling
law.  The point \(p\) belongs to at most one popular circle.  If there is
one, it differs from the unit circle centered at \(q\), since the
latter is not popular.  Their intersection has at most two points.
In this original law, \(r\) is a uniform neighbor of \(q\), so the
probability that it shares a popular circle with \(p\) is at most
\(2/d(q)\), even before requiring \(r\notin P_{\rm mult}\).
This proves \eqref{ext:circle-cost}.

For the comparison with the selected band, recall
\(R_0=k^{1/8}\), \(k=t^{o(1)}\), and the explicit requirement
\(k^{20}/\min_qd(q)\to0\) in \Cref{bud:parameters}.
The first two terms of \eqref{ext:circle-cost}, multiplied by
\(R_0\log(2R_0)\), tend to zero by a fixed power of \(t\).
For the third term,
\[
 \frac{R_0\log(2R_0)}{\min_qd(q)}
 =\frac{k^{1/8}\log(2k^{1/8})}{k^{20}}
                  \frac{k^{20}}{\min_qd(q)}\longrightarrow0.
\]
This establishes the required relative loss bound.
\end{proof}

We can now forget the probability weights and retain an ordinary graph.
The estimates already hold separately for every surviving wedge, so
choosing one center for each endpoint pair preserves them.

\begin{proposition}[Organized pair graph]\label{ext:pairs}
For the incidence graphs and levels constructed above, there are
\(J\to\infty\), a sequence \(\omega_t\to0\), and a simple bipartite graph
\(\mathcal E\) between two copies of \(P\), each of size at most
\(t^{1+o(1)}\), with at least \(t^{2-o(1)}\) edges, having the following
properties.  Each edge \((p,r)\) has distinct endpoints and an assigned
center \(q=q(p,r)\in Q\) with \(pq,rq\in E\).  Set
\[
 \lambda_{pr}=\frac{z_p-z_{q(p,r)}}{z_r-z_{q(p,r)}}.
\]
No edge has both endpoints on any real unit circle containing at least
\(t^{9/10}\) points of \(P\).  Uniformly over the edges and complete
rectangles of \(\mathcal E\),
\begin{align}
 (z_p-z_r)(w_p-w_r)&=2-\lambda_{pr}-\lambda_{pr}^{-1},
                                               \label{ext:pair-identity}\\
 h(\lambda_{pr})&\ge J/3,                      \label{ext:pair-height}\\
 h\!\left(\frac{\lambda_{pr}\lambda_{p'r'}}
                    {\lambda_{pr'}\lambda_{p'r}}\right)
 &\le\omega_tJ.                               \label{ext:rectangle-height}
\end{align}
In the last line all four indicated edges are required to be present;
their assigned centers need not agree.  The edge ratios
\(\lambda_{pr}\) are unrelated to the fixed score constant \(\lambda\).
\end{proposition}

\begin{proof}
Take the band from \Cref{ext:band}, and remove the wedges excluded by
\Cref{ext:circles}.  Its probability remains at least
\(c/(R_0\log(2R_0))=t^{-o(1)}\), and all its endpoint pairs are distinct.
For any fixed off-diagonal pair \((p,r)\), its ordinary wedge probability is
\[
 \sum_{q:\,pq,rq\in E}\frac{1}{\abs E\,d(q)}
 \le\frac{2}{\abs E\min_qd(q)}.
\]
Here the sum has at most two terms, since two distinct real points have
at most two unit-circle centers.  The retained wedges therefore contain
at least
\[
 \frac{c\abs E\min_qd(q)}{2R_0\log(2R_0)}\ge t^{2-o(1)}
\]
distinct ordered endpoint pairs, using
\(\abs E=t^{2+2a+o(a)}\) and \(\min_qd(q)=t^{a+o(a)}\).
These pairs form \(\mathcal E\).  Choose one surviving witnessing
center for each pair.

For its two jumps \(d=z_p-z_q\), \(d'=z_r-z_q\), the exact unit equations
give \(w_p-w_q=1/d\) and \(w_r-w_q=1/d'\).  Hence
\[
 (z_p-z_r)(w_p-w_r)=(d-d')(1/d-1/d')
                    =2-d/d'-d'/d,
\]
proving \eqref{ext:pair-identity}.  The product formula gives
\[
 2h(\lambda_{pr})=\sum_v^*\abs{\log\abs{\lambda_{pr}}_v}
                  =\norm{n_{pq}-n_{rq}}_1
                  \ge(1-\delta_t)J,
\]
so \eqref{ext:pair-height} holds for all sufficiently large indices.

Finally, put
\[
 b_{pr}(v)=\log\abs{\lambda_{pr}}_v+a_p(v)-a_r(v).
\]
By \eqref{ext:profile-error} and \eqref{ext:band-conditions},
\(\norm{b_{pr}}_1\le2\delta_tJ\), uniformly on all retained edges.
On a complete rectangle, the unary profiles cancel, leaving
\[
 \log\left|\frac{\lambda_{pr}\lambda_{p'r'}}
                    {\lambda_{pr'}\lambda_{p'r}}\right|_v
     =b_{pr}(v)+b_{p'r'}(v)-b_{pr'}(v)-b_{p'r}(v).
\]
Its weighted absolute sum is at most \(8\delta_tJ\).  Another use of the
product formula bounds its height by \(4\delta_tJ\), proving
\eqref{ext:rectangle-height} with \(\omega_t=4\delta_t\).
The popular-circle exclusion holds by construction.
\end{proof}
